\section{Thông tin nhóm}

\subsection{Danh sách thành viên}
\label{subsec:danh-sach-thanh-vien}

Nhóm \textbf{RushDL - 13}, lớp \textbf{24KTPM2}, gồm hai thành viên:

\begin{center}
\captionof{table}{Danh sách thành viên nhóm}
\label{tab:thanh-vien}
\begin{tabularx}{0.9\textwidth}{|C{1cm}|Y|C{3.3cm}|}
  \hline
  \textbf{STT} & \textbf{Họ và tên} & \textbf{MSSV}\\
  \hline
  1 & Nguyễn Sơn Hải & 24127163\\
  \hline
  2 & Châu Vũ Trung & 24127256\\
  \hline
\end{tabularx}
\end{center}

\subsection{Mục tiêu làm việc của nhóm}

Nhóm thực hiện đồ án Hệ thống chat nhằm vận dụng kiến thức của học phần
Lập trình ứng dụng Java vào việc xây dựng một ứng dụng có tổ chức, đáp ứng
các yêu cầu về quản lý người dùng và trao đổi tin nhắn. Thông qua đồ án,
các thành viên hướng đến việc rèn luyện khả năng phân tích yêu cầu, thiết kế
ứng dụng và phối hợp phát triển phần mềm.

Trong Version 1 (20\%), nhóm tập trung xác định phạm vi hệ thống, thiết kế giao diện
cho hai phân hệ quản trị viên và người sử dụng, đồng thời thiết kế cơ sở dữ liệu
phù hợp với các chức năng trong đề bài. Các sản phẩm của giai đoạn này là
nền tảng để triển khai và hoàn thiện chức năng ở phiên bản tiếp theo.
Theo yêu cầu của giai đoạn này, các màn hình chưa cần lập trình xử lý sự kiện.

Về tổ chức làm việc, nhóm hướng đến việc phân chia trách nhiệm rõ ràng,
hoàn thành công việc đúng thời hạn và bảo đảm nội dung đáp ứng đầy đủ yêu cầu
của đề bài. Nhóm sử dụng Messenger để trao đổi và cập nhật tiến độ;
Google Drive để lưu trữ, chia sẻ tài liệu và phối hợp hoàn thiện báo cáo;
Git và GitHub để quản lý mã nguồn, theo dõi thay đổi và ghi nhận đóng góp
của từng thành viên.

\noindent Các kênh lưu trữ và quản lý sản phẩm của nhóm:
\begin{itemize}
  \item \textbf{Tài liệu và báo cáo:}
    \href{https://drive.google.com/drive/folders/1Kci4Ug_HXLRbBBFc0h6lVz0RbuvbcJTF?usp=sharing}
    {Thư mục Google Drive của nhóm}.
  \item \textbf{Mã nguồn:}
    \href{https://github.com/nguyenshai/Project-java}
    {Repository GitHub của nhóm}.
\end{itemize}

\clearpage
\subsection{Vai trò và trách nhiệm của các thành viên}
\label{subsec:vai-tro-thanh-vien}

Nhóm dự kiến phân chia trách nhiệm theo hai phân hệ. Nguyễn Sơn Hải tập trung
vào các chức năng quản trị và theo dõi hệ thống; Châu Vũ Trung tập trung vào
các chức năng phục vụ người dùng và quá trình giao tiếp. Cách phân chia này
giúp xác định rõ phạm vi phụ trách và hạn chế chồng chéo công việc.

\begin{center}
\small
\captionof{table}{Phạm vi phụ trách của từng thành viên}
\label{tab:vai-tro-thanh-vien}
\begin{tabularx}{\textwidth}{|L{3cm}|L{2.8cm}|Y|}
  \hline
  \textbf{Thành viên} & \textbf{Phạm vi phụ trách} &
  \textbf{Trách nhiệm chính trong V1}\\
  \hline
  Nguyễn Sơn Hải\newline 24127163 &
  Phân hệ quản trị viên &
  Phụ trách thiết kế giao diện phân hệ Admin; chuẩn bị nội dung báo cáo
  và ảnh minh họa tương ứng.\\
  \hline
  Châu Vũ Trung\newline 24127256 &
  Phân hệ người sử dụng &
  Phụ trách thiết kế giao diện phân hệ User; chuẩn bị nội dung báo cáo
  và ảnh minh họa tương ứng; tổng hợp báo cáo V1.\\
  \hline
\end{tabularx}
\end{center}

Hai thành viên cùng tham gia thiết kế cơ sở dữ liệu, thống nhất cách tổ chức
dự án và kiểm tra chéo sản phẩm. Mỗi người chịu trách nhiệm cung cấp nội dung
chính xác cho phần mình phụ trách; người tổng hợp báo cáo có trách nhiệm
thống nhất cách trình bày và kiểm tra sự đầy đủ của các mục.

Phân công trên là đề xuất ban đầu, có thể điều chỉnh theo khối lượng thực tế
sau khi hai thành viên thống nhất. Các đầu việc và phạm vi chức năng liên quan
được cụ thể hóa tại mục~\ref{subsec:ke-hoach-phan-cong}.

\clearpage
\subsection{Quy tắc phối hợp và thực hiện công việc}

Nhóm đề xuất áp dụng các quy tắc sau:

\begin{enumerate}
  \item \textbf{Bám sát yêu cầu:}
    Trước khi thực hiện, mỗi thành viên cần xác định rõ nội dung và sản phẩm
    cần hoàn thành. Phần việc phải đáp ứng các yêu cầu tương ứng của đề bài,
    không tự ý bỏ sót hoặc thay đổi phạm vi.

  \item \textbf{Tuân thủ thời hạn:}
    Mỗi đầu việc cần có thời hạn hoàn thành được nhóm thống nhất, phù hợp
    với hạn nộp của học phần và có thời gian dành cho kiểm tra, chỉnh sửa.

  \item \textbf{Chủ động báo cáo tiến độ:}
    Các thành viên cập nhật tình hình qua Messenger. Khi có khó khăn hoặc
    nguy cơ trễ hạn, cần thông báo sớm để nhóm điều chỉnh hoặc hỗ trợ.

  \item \textbf{Quản lý sản phẩm chung:}
    Tài liệu được chia sẻ trên Google Drive; mã nguồn được quản lý bằng
    Git và GitHub. Mỗi thành viên thực hiện commit cho phần mình đóng góp,
    với nội dung mô tả rõ thay đổi.

  \item \textbf{Kiểm tra chéo:}
    Sau khi hoàn thành phần việc, thành viên còn lại rà soát tính đầy đủ,
    sự thống nhất và mức độ đáp ứng yêu cầu trước khi tổng hợp.

  \item \textbf{Rà soát trước khi nộp:}
    Cả hai cùng kiểm tra sản phẩm, nội dung báo cáo, cấu trúc hồ sơ và
    minh chứng commit, bảo đảm hồ sơ đầy đủ và được nộp đúng hạn.

  \item \textbf{Ghi nhận đóng góp thực tế:}
    Phân công và tỷ lệ đóng góp được cập nhật theo công việc thực sự đã
    thực hiện, có sự thống nhất của cả hai thành viên.
\end{enumerate}

\noindent\textbf{Nguyên tắc nghiệm thu:}
Mọi công việc đều phải được cả hai thành viên kiểm tra và thống nhất
công nhận trước khi ghi nhận hoàn thành.
